\chapter{Famílies, fibres i operacions indexades}\label{cap:families}

\begin{objectius}
\apartat{Prerequisits} Del capítol~\ref{cap:parts}: pullbacks, imatge inversa i fibres. Del capítol~\ref{cap:relacions}: seccions, l'axioma d'elecció i les particions. Del capítol~\ref{cap:parts-interna}: conjunt de parts, pertinença interna i els quantificadors \(\exists_f\), \(\forall_f\) amb les seves lectures i adjuncions. Dels capítols~\ref{cap:infinitud} i~\ref{cap:cantor}: les famílies numerables i la construcció d'ínfims arbitraris amb \(\forall\).
\apartat{Objectius} En acabar el capítol, el lector hauria de poder:
\begin{itemize}
  \item representar una família de conjunts indexada per \(I\) com una aplicació \(p\colon E\to I\), i llegir-ne les fibres com a parts del total;
  \item identificar seccions amb funcions d'elecció, i situar l'elecció com l'existència de seccions per a famílies de fibres no buides;
  \item construir la suma (el total) i el producte (les seccions, dins \(E^I\)) d'una família, i distingir l'existència del producte com a objecte de l'existència dels seus punts;
  \item codificar una família de parts d'un conjunt fix com \(\alpha\colon I\to P(A)\), construir-ne la unió i la intersecció indexades amb els quantificadors, i demostrar-ne les propietats de suprem, ínfim i De Morgan;
  \item demostrar que una família de parts disjuntes que cobreixen \(A\) és una partició, tancant el pont del capítol~\ref{cap:relacions};
  \item distingir família interna de regla externa, aplicar la regla de treball, i saber què promet l'esquema SR i per què no és l'elecció.
\end{itemize}
\end{objectius}

Fins ara hem indexat col·leccions de conjunts d'una manera implícita: una
família \((A_n)_{n\in\N}\) apareixia al capítol~\ref{cap:infinitud} com una
aplicació \(p\colon E\to\N\) amb les fibres per components
(definició~\ref{def:familia-numerable}), i els arguments diagonals del
capítol~\ref{cap:cantor} només comparaven conjunts d'un en un. Els capítols
transfinits que vindran ---recursió sobre un bon ordre, suprems de tipus
ordinals, successions d'alephs--- necessiten manejar \emph{famílies} de
conjunts amb tota generalitat, i sobretot necessiten un criteri net per
distingir dues coses que en la pràctica clàssica es confonen: una família
que \emph{ja existeix} com a objecte, i una \emph{regla} que a cada índex
assigna un conjunt però que encara no s'ha reunit en cap objecte. Aquesta
distinció és el cor del capítol i la clau de tot el que segueix; ja es va
insinuar amb l'assignació \(n\mapsto\seg{n}\) del capítol~\ref{cap:naturals},
i reapareixerà a cada recursió transfinita.

La idea central és simple i no demana cap vocabulari nou. Una família de
conjunts indexada per un conjunt \(I\) es representa amb una aplicació
\(p\colon E\to I\); el conjunt \(E\) és el \emph{total}, i el conjunt indexat
per \(i\) és la \emph{fibra} de \(p\) sobre \(i\). Tot el capítol desplega
aquesta idea ---suma, producte, unions i interseccions indexades--- amb les
eines que ja tenim: pullbacks, exponencials, el conjunt de parts i els
quantificadors interns. No parlarem de functors ni de fibracions; no en
calen.

\section{Famílies com a aplicacions, i les seves fibres}\label{sec:families-aplicacions}

\begin{definicio}[Família de conjunts]\index{Família de conjunts}\label{def:familia}
Una \emph{família de conjunts indexada per \(I\)} és una aplicació
\(p\colon E\to I\). El conjunt \(I\) és l'\emph{índex}, \(E\) és el conjunt
\emph{total}, i per a cada element \(i\colon\mathbf1\to I\) la \emph{fibra}
de la família sobre \(i\) és el producte fibrat
\[
E_i\ :=\ \mathbf1\times_I E
\]
de \(i\colon\mathbf1\to I\) i \(p\colon E\to I\)
(definició~\ref{def:pullback}), amb la projecció
\(E_i\hookrightarrow E\) que n'és un monomorfisme
(proposició~\ref{prop:pullback-mono}). Escrivim també \(A_i:=E_i\) quan
volem pensar la família com \((A_i)_{i\in I}\).
\end{definicio}

\begin{observacio}[La fibra és una part del total]\label{obs:fibra-part}
La fibra \(E_i=\mathbf1\times_I E\) és exactament la imatge inversa
\(p^{-1}(i)\hookrightarrow E\) del punt \(i\) per \(p\)
(definició~\ref{def:imatge-inversa}, i
observació~\ref{obs:pullback-exemples}): en la lectura clàssica, els seus
punts són els \(e\in E\) amb \(p(e)=i\). Per tant una família \(p\colon E\to
I\) no és res més que el seu total \(E\) \emph{repartit} en parts disjuntes,
una per cada índex, i cada part etiquetada pel seu índex a través de \(p\).
Aquesta és la manera estructural de dir \enquote{una família de conjunts}
sense sortir dels conjunts ja construïts: tot viu dins \(E\), i les fibres
són parts seves.
\end{observacio}

L'exemple més senzill és el d'una família que no varia amb l'índex.

\begin{definicio}[Família constant]\index{Família constant}\label{def:familia-constant}
Donat un conjunt \(A\), la \emph{família constant} de valor \(A\) indexada
per \(I\) és la projecció
\[
\pi_I\colon I\times A\longrightarrow I .
\]
La seva fibra sobre cada \(i\) és \(\{i\}\times A\cong A\)
(observació~\ref{obs:pullback-exemples}), de manera que representa la família
\((A)_{i\in I}\) amb totes les fibres iguals a \(A\).
\end{definicio}

\begin{observacio}[Tota aplicació és una família, i tota part una família de valors booleans]\label{obs:familia-exemples}
La noció és prou àmplia per recobrir casos ja coneguts. Una aplicació
qualsevol \(p\colon E\to I\) es \emph{pot} llegir com a família, i triar com
llegir-la depèn de què vulguem: la mateixa \(p\colon E\to\N\) del
capítol~\ref{cap:infinitud} era la família numerable
(definició~\ref{def:familia-numerable}), cas particular amb \(I=\N\). Una
part \(m\colon S\hookrightarrow I\), vista com la seva característica
\(\chi_S\colon I\to\mathbf2\), és la família indexada per \(I\) amb fibra
\(\mathbf1\) sobre els \(i\in S\) i \(\mathbf0\) fora: una família de conjunts
\enquote{de com a molt un element}. Aquests exemples fixen la intuïció; el
contingut del capítol és el que se'n construeix.
\end{observacio}

\section{Seccions i funcions d'elecció}\label{sec:seccions}

Una família \(p\colon E\to I\) reparteix \(E\) en fibres; \emph{triar} un
element a cada fibra és exactament donar una secció de \(p\).

\begin{definicio}[Secció d'una família]\index{Secció d'una família}\label{def:seccio-familia}
Una \emph{secció} de la família \(p\colon E\to I\) és una aplicació
\(s\colon I\to E\) amb \(p\circ s=\id_I\). Per a cada índex \(i\), el valor
\(s(i)\) és un element de la fibra \(E_i\), perquè \(p(s(i))=i\); una secció
és, doncs, una \emph{tria simultània} d'un element a cada fibra ---una
\emph{funció d'elecció} per a la família.
\end{definicio}

\begin{observacio}[Quan existeix una secció: el lloc exacte de l'elecció]\label{obs:seccio-eleccio}
L'existència d'una secció \emph{no} és automàtica, i és precisament aquí on
apareix ---o no--- l'axioma d'elecció. (Diem indistintament que una fibra és
\emph{no inicial} o \emph{no buida}: per la separació global i el lema del
buit, un conjunt és inicial si i només si no té cap punt global, de manera
que les dues expressions coincideixen en aquest llibre.) Si totes les fibres són no inicials,
és a dir si \(p\) és un epimorfisme (exhaustivitat sobre punts,
capítol~\ref{cap:valors}; i tot epimorfisme és regular,
capítol~\ref{cap:imatges}), l'axioma d'elecció
(axioma~\ref{ax:eleccio}) diu exactament que \(p\) admet una secció
(proposició~\ref{prop:epi-regular-eleccio}). Sense elecció, en canvi, una
família de fibres totes no buides pot no tenir cap secció. Convé retenir la
distinció, que serà constant en els capítols transfinits: \emph{tenir
fibres no buides} és una propietat de la família ja donada; \emph{tenir una
secció} és una tria simultània que, en general, costa l'elecció. Quan les
fibres porten un element distingit canònic ---un mínim, un ínfim--- la secció
es construeix sense elecció, com vam veure amb els subobjectes de \(\N\)
(observació~\ref{obs:minim-vs-eleccio}) i amb la part \(D\) de
Schröder--Bernstein (observació~\ref{obs:sb-sense-eleccio}).
\end{observacio}

\section{Suma i producte d'una família}\label{sec:suma-producte}

Les dues operacions bàsiques sobre una família són reunir totes les fibres
en un sol conjunt ---la suma--- i formar el conjunt de totes les seccions
---el producte. Totes dues es construeixen amb el que ja tenim.

\begin{definicio}[Suma d'una família]\index{Suma d'una família}\label{def:suma-familia}
La \emph{suma} de la família \(p\colon E\to I\) és el seu propi total \(E\),
pensat com la reunió disjunta de les fibres,
\[
\sum_{i\in I}E_i\ :=\ E,
\]
en què cada fibra \(E_i\hookrightarrow E\) apareix marcada pel seu índex a
través de \(p\). L'aplicació \(p\) mateixa recorda, per a cada element del
total, a quina fibra pertany.
\end{definicio}

\begin{observacio}[La suma no demana res de nou]\label{obs:suma-existeix}
La suma d'una família \emph{existeix sempre}, sense cap axioma ni cap tria:
és el conjunt \(E\), que és una de les dades. Aquesta trivialitat aparent és
important. Recollir \enquote{totes les fibres alhora} és exactament tenir el
total \(E\); el que pot mancar, com veurem, no és la reunió sinó la
\emph{regla} que produeix \(E\) a partir d'una assignació índex-a-conjunt
quan aquesta assignació no ve ja d'una aplicació. Per a la família constant
\(I\times A\to I\), la suma és \(I\times A\), i recupera el producte
cartesià habitual llegit com \enquote{\(I\) còpies de \(A\)}.
\end{observacio}

\begin{definicio}[Producte d'una família]\index{Producte d'una família}\label{def:producte-familia}
El \emph{producte} de la família \(p\colon E\to I\) és el conjunt de les
seves seccions,
\[
\prod_{i\in I}E_i\ :=\ \{\,s\in E^{I}\ \mid\ p\circ s=\id_I\,\}
\ \hookrightarrow\ E^{I},
\]
la part de l'exponencial \(E^{I}\) (axioma~\ref{ax:exponencials}) tallada per
l'equació \(p\circ s=\id_I\) ---l'equalitzador de les dues aplicacions
paral·leles \(E^{I}\to I^{I}\) donades per \(s\mapsto p\circ s\)
(postcomposició interna amb \(p\)) i per la constant \(s\mapsto\id_I\). Els seus elements
globals són exactament les seccions de \(p\) (definició~\ref{def:seccio-familia}).
\end{definicio}

\begin{proposicio}[Existència del producte i el lloc de l'elecció]\label{prop:producte-existeix}
Per a tota família \(p\colon E\to I\), el producte \(\prod_{i\in I}E_i\)
\emph{existeix com a conjunt}: és una part interna ben definida de
\(E^{I}\). En canvi, que sigui \emph{no inicial} ---que tingui algun element
global--- equival a l'existència d'una secció de \(p\), i per tant, quan les
fibres són no buides, a una forma d'elecció.\par
\emph{Axiomes usats:} exponencials (axioma~\ref{ax:exponencials}),
equalitzadors i parts (capítols~\ref{cap:parts} i~\ref{cap:valors}); cap
axioma nou per a l'existència \emph{com a objecte}. La no-vacuïtat, per a
fibres no buides, usa l'elecció (axioma~\ref{ax:eleccio}).
\end{proposicio}

\begin{prova}
L'exponencial \(E^{I}\) existeix per l'axioma~\ref{ax:exponencials}, i les
dues aplicacions paral·leles \(E^{I}\to I^{I}\) donades per
\(s\mapsto p\circ s\) (postcomposició interna amb \(p\): la composició
interna del capítol~\ref{cap:exponencials} amb el primer argument fixat al
nom de \(p\)) i per \(s\mapsto\ulcorner\id_I\urcorner\) (constant) tenen un
equalitzador, que és una part de \(E^{I}\) (capítol~\ref{cap:parts}); aquesta part és, per construcció,
\(\prod_{i\in I}E_i\). Fins aquí, cap tria. Un element global del producte és
una \(s\colon\mathbf1\to E^{I}\) amb \(p\circ s=\id_I\), és a dir ---per
transposició (definició~\ref{def:exponencial})--- una aplicació
\(I\to E\) amb \(p\circ s=\id_I\): una secció. Per tant el producte és no
inicial si i només si \(p\) té una secció. Si les fibres són totes no buides
(\(p\) epimorfisme), l'existència d'una tal secció és exactament el
que afirma l'axioma d'elecció (axioma~\ref{ax:eleccio});
sense elecció pot fallar.
\end{prova}

\begin{observacio}[Objecte contra elements: una distinció que cal mantenir]\label{obs:objecte-vs-elements}
La proposició anterior separa dues qüestions que la notació
\(\prod_{i\in I}E_i\) tendeix a confondre: \emph{existeix el producte?} i
\emph{té elements?}. La primera es respon sempre que sí, amb exponencials i
equalitzadors. La segona, per a fibres no buides, és el contingut mateix de
l'elecció. Aquesta és la versió general del que vam veure amb el producte
numerable (observació~\ref{obs:producte-numerable}): l'objecte hi és sempre;
els seus punts, no necessàriament. La mateixa distinció ---objecte contra
punts globals--- reapareixerà amb les famílies definides per regles, a la
secció~\ref{sec:regla-vs-familia}.
\end{observacio}

\section{Famílies de parts d'un conjunt fix}\label{sec:families-parts}

Un cas particularment manejable, i el que més usarem, és el de les famílies
totes les fibres de les quals són parts d'un mateix conjunt fix \(A\).
Aleshores la família es codifica com una aplicació cap al conjunt de parts, i
totes les operacions es fan \emph{acotadament} dins \(A\) i \(P(A)\).

\begin{definicio}[Família de parts d'un conjunt fix]\index{Família de parts d'un conjunt fix}\label{def:familia-parts}
Una \emph{família de parts} de \(A\) indexada per \(I\) és una aplicació
\[
\alpha\colon I\longrightarrow P(A),
\]
que a cada índex \(i\) assigna una part \(\alpha(i)\hookrightarrow A\)
(el seu nom, definició~\ref{def:conjunt-parts}). La \emph{relació de
pertinença total} de la família és la part
\[
R\ :=\ \{\,(i,a)\in I\times A\ \mid\ a\mathrel{\varepsilon}\alpha(i)\,\}
\ \hookrightarrow\ I\times A,
\]
la imatge inversa de la pertinença interna \({\in_A}\)
(definició~\ref{def:pertinenca-interna}) per \(\id_I\times\alpha\), seguida
de la commutació. Escrivim \(R_i\hookrightarrow A\) per a la fibra de la
projecció \(R\to I\) sobre \(i\), que és exactament \(\alpha(i)\).
\end{definicio}

\begin{observacio}[Tot és acotat]\label{obs:familia-parts-acotada}
Aquesta és la situació ideal per als capítols transfinits, i la raó per la
qual la guia insisteix a \emph{acotar} les col·leccions. Una família de parts
de \(A\) és un únic objecte, \(\alpha\colon I\to P(A)\), que viu dins un
exponencial fix \(P(A)^{I}\); la seva relació de pertinença total \(R\) viu
dins \(I\times A\); i cada fibra és una part de \(A\). Res no s'escapa a un
ambient prèviament construït. Sempre que puguem expressar una col·lecció de
conjunts com a família de parts d'un fix, les seves unions, interseccions i
suprems es formaran sense cap principi nou, com veurem tot seguit. La
dificultat transfinita començarà exactament quan una col·lecció \emph{no} es
pugui allotjar així.
\end{observacio}

\begin{definicio}[Unió i intersecció indexades]\index{Unió i intersecció indexades}\label{def:unio-interseccio-indexada}
Sigui \(\alpha\colon I\to P(A)\) una família de parts de \(A\), amb relació
de pertinença total \(R\hookrightarrow I\times A\) i projecció
\(\pi_A\colon I\times A\to A\). La \emph{unió indexada} i la
\emph{intersecció indexada} de la família són les parts de \(A\)
\[
\bigcup_{i\in I}\alpha(i)\ :=\ \exists_{\pi_A}(R),
\qquad
\bigcap_{i\in I}\alpha(i)\ :=\ \forall_{\pi_A}(R),
\]
on \(\exists_{\pi_A}\) i \(\forall_{\pi_A}\) són els quantificadors interns
associats a la projecció (definició~\ref{def:quantificadors}), aplicats al
nom de la part \(R\). Sobre punts globals (proposició~\ref{prop:quantificadors-lectura}):
un punt \(a\) pertany a la unió quan existeix un punt \(i\) amb
\(a\mathrel{\varepsilon}\alpha(i)\), i a la intersecció quan
\(a\mathrel{\varepsilon}\alpha(i)\) per a tot \(i\); sobre elements
generalitzats, la lectura de la unió és amb recobriment, com sempre amb
\(\exists\).
\end{definicio}

\begin{observacio}[La lectura amb quantificadors]\label{obs:unio-interseccio-lectura}
La formulació anterior és la traducció directa de \enquote{existeix un
índex\ldots} i \enquote{per a tot índex\ldots} al llenguatge intern del
capítol~\ref{cap:parts-interna}. Per a la unió: \(a\in\bigcup_i\alpha(i)\)
diu que la fibra de \(R\) sobre \(a\) ---els índexs \(i\) amb
\(a\mathrel{\varepsilon}\alpha(i)\)--- és no buida, i això és exactament
\(\exists_{\pi_A}(R)\). Per a la intersecció, \(a\in\bigcap_i\alpha(i)\) diu
que \emph{tot} índex compleix \(a\mathrel{\varepsilon}\alpha(i)\), és a dir
\(\forall_{\pi_A}\) del predicat \enquote{\((i,a)\in R\)}: tots els punts de
la fibra de \(\pi_A\) sobre \(a\), que són els \((i,a)\), han de ser a \(R\).
La construcció usa \emph{només} els
quantificadors interns, ja disponibles; cap suprem infinit adjuntat, cap
axioma nou. És la mateixa tècnica amb què vam formar la part \(D\) de
Schröder--Bernstein com un ínfim arbitrari
(teorema~\ref{teo:schroder-bernstein}).
\end{observacio}

\begin{proposicio}[Propietats de la unió i la intersecció indexades]\label{prop:unio-interseccio-propietats}
Sigui \(\alpha\colon I\to P(A)\) una família de parts de \(A\). Aleshores,
dins \(P(A)\):
\begin{enumerate}
\item \(\bigcup_{i\in I}\alpha(i)\) és la menor part de \(A\) que conté totes
les \(\alpha(i)\); és a dir, el suprem de la família a \((P(A),\subseteq)\).
\item \(\bigcap_{i\in I}\alpha(i)\) és la major part de \(A\) continguda en
totes les \(\alpha(i)\); és a dir, l'ínfim de la família.
\item Es compleixen les lleis de De Morgan indexades:
\[
A\setminus\bigcup_{i\in I}\alpha(i)=\bigcap_{i\in I}\bigl(A\setminus\alpha(i)\bigr),
\qquad
A\setminus\bigcap_{i\in I}\alpha(i)=\bigcup_{i\in I}\bigl(A\setminus\alpha(i)\bigr).
\]
\end{enumerate}
\emph{Axiomes usats:} conjunt de parts, pertinença interna i quantificadors
(capítol~\ref{cap:parts-interna}), complement (capítol~\ref{cap:valors}); cap
axioma nou.
\end{proposicio}

\begin{prova}
(1) i (2) són les propietats universals de \(\exists_{\pi_A}\) i
\(\forall_{\pi_A}\) com a adjunts de la imatge inversa
(proposició~\ref{prop:quantificadors-adjuncio}, amb l'estadi \(X=\mathbf1\)):
\(\exists_{\pi_A}(R)\subseteq S\) equival a \(R\subseteq\pi_A^{-1}(S)=I\times S\),
és a dir, sobre punts, a \(\alpha(i)\subseteq S\) per a tot \(i\) ---la
propietat del suprem---; i dualment, \(S\subseteq\forall_{\pi_A}(R)\) equival
a \(I\times S\subseteq R\), és a dir, a \(S\subseteq\alpha(i)\) per a tot
\(i\) ---la de l'ínfim. (3) surt de les mateixes adjuncions per
complementació, com a la construcció de \(\forall_f=\neg\exists_f\neg\) de la
proposició~\ref{prop:quantificadors-adjuncio}, transportant les lleis de De
Morgan booleanes (observació~\ref{obs:lleis-transportades}) a la família.
L'argument no fa servir cap hipòtesi sobre \(I\): si \(I\) és inicial, la
unió és la part buida i la intersecció és tot \(A\), que són, en efecte, el
suprem i l'ínfim de la família buida.
\end{prova}

Amb aquest llenguatge podem tancar el pont que el capítol~\ref{cap:relacions}
va deixar obert: una partició estructural (un epimorfisme) té fibres
disjuntes que cobreixen el conjunt (proposició~\ref{prop:fibres-particio});
recíprocament, una família de parts amb aquestes propietats és una
partició.

\begin{proposicio}[Les famílies de parts disjuntes que cobreixen són particions]\label{prop:familia-particio}
Sigui \(\alpha\colon I\to P(A)\) una família de parts de \(A\) tal que,
sobre punts: cada \(\alpha(i)\) és no buida; dos índexs diferents donen
parts disjuntes, \(i\ne j\Rightarrow\alpha(i)\cap\alpha(j)=\varnothing_A\);
i la unió indexada és tot \(A\). Aleshores existeix un epimorfisme
\(q\colon A\twoheadrightarrow I\) ---una partició--- les fibres del qual són
exactament les \(\alpha(i)\), i és únic amb aquesta propietat.\par
\emph{Axiomes usats:} els del capítol~\ref{cap:parts-interna} i la separació
global; cap elecció.
\end{proposicio}

\begin{prova}
Sigui \(R\hookrightarrow I\times A\) la relació de pertinença total de la
família, amb projeccions \(r_I\colon R\to I\) i \(r_A\colon R\to A\). La
hipòtesi que les parts cobreixen \(A\) diu que \(\exists_{\pi_A}(R)\) és
total, és a dir, que \(r_A\) és exhaustiva sobre punts: tot punt \(a\)
pertany a alguna \(\alpha(i)\). La disjunció diu que \(r_A\) és injectiva
sobre punts: si \((i,a)\) i \((j,a)\) són punts de \(R\), aleshores
\(a\in\alpha(i)\cap\alpha(j)\), que és buida si \(i\ne j\); per tant \(i=j\).
Per la separació global, \(r_A\) és un bimorfisme i, doncs, un isomorfisme
(proposició~\ref{prop:bimorfisme-isomorfisme}). Posem
\(q:=r_I\circ r_A^{-1}\colon A\to I\). És exhaustiva sobre punts, perquè
cada \(\alpha(i)\) té un punt \(a\), i aleshores \((i,a)\in R\) i
\(q(a)=i\); per tant és un epimorfisme, i una partició. La fibra
\(q^{-1}(i)\) és, per construcció, la part dels \(a\) amb \((i,a)\in R\), és
a dir, \(\alpha(i)\). La unicitat: una altra \(q'\) amb les mateixes fibres
coincideix amb \(q\) sobre tots els punts, i per la separació global
\(q'=q\).
\end{prova}

\section{Regla externa contra família interna}\label{sec:regla-vs-familia}

Arribem al punt central del capítol, la distinció que la guia demana fer molt
explícita perquè governa tota la teoria transfinita. Cal separar dos usos de
la notació \enquote{\(i\mapsto X_i\)}.

\begin{definicio}[Família interna i regla externa]\index{Família interna i regla externa}\label{def:familia-vs-regla}
Una \emph{família interna} indexada per \(I\) és una aplicació \emph{ja
donada} \(p\colon E\to I\) (definició~\ref{def:familia}); les seves fibres
\(E_i\) són parts d'un total \(E\) que existeix com a objecte. En canvi, una
\emph{regla externa} és una prescripció que, al metallenguatge, associa a
cada element \(i\) d'\(I\) un conjunt \(X_i\), \emph{sense} que se n'hagi
exhibit cap aplicació \(E\to I\) les fibres de la qual siguin els \(X_i\).
Una regla externa \emph{es representa} per una família interna quan existeix
una tal \(p\colon E\to I\) amb \(E_i\cong X_i\) per a tot \(i\), de manera
compatible.
\end{definicio}

\begin{observacio}[El problema, i on ja el vam trobar]\label{obs:regla-externa-problema}
La diferència no és pedant: hi ha regles externes perfectament definides que
\emph{no} es poden representar per cap família interna amb els axiomes que
tenim. L'exemple prototípic ja va aparèixer al capítol~\ref{cap:naturals}: la
regla \(n\mapsto\seg{n}\), que a cada natural assigna el seu segment inicial,
és una prescripció clara, però reunir tots els \(\seg{n}\) en un sol total
\(E\to\N\) demana una construcció que no és cap dels axiomes bàsics ---de fet
es fa allotjant tots els \(\seg{n}\) dins un exponencial fix, com al
conjunt de llistes (definició~\ref{def:llistes}). Sempre que puguem allotjar
els \(X_i\) dins un conjunt ambient comú ---per exemple si tots són parts
d'un fix \(A\), i aleshores la regla és una aplicació \(I\to P(A)\)
(definició~\ref{def:familia-parts})--- la representació existeix i no cal res
nou. El problema genuí apareix quan els \(X_i\) \emph{creixen sense cota
comuna coneguda}: aleshores no hi ha cap ambient dins el qual tallar-los per
separació, i la reunió pot escapar-se dels axiomes.
\end{observacio}

\begin{observacio}[La regla de treball]\label{obs:regla-treball}
D'aquí surt una disciplina que mantindrem a tot el desenvolupament
transfinit: \textbf{cada vegada que aparegui una notació \(i\mapsto X_i\),
cal indicar quina aplicació \(E\to I\) representa la família; i si no
existeix encara, cal declarar-ho.} Escriure \(\bigcup_{i\in I}X_i\) o
\(\sup_i X_i\) només té sentit un cop s'ha exhibit el total \(E\), o un
ambient \(A\) que contingui tots els \(X_i\). Aquesta és la lliçó que fa
manejable la recursió transfinita: mentre els valors visquin dins un
codomini fix, la família és interna i tot és de Nivell~I; quan els valors
siguin conjunts nous sense codomini comú, cal un principi addicional, que
introduïm a la secció següent i no abans.
\end{observacio}

\section{Nota sobre la reunió de famílies no acotades (SR)}\label{sec:sr}

\begin{nota}
Aquesta secció és de \emph{Nivell~II}: enuncia, però \emph{no adopta},
un principi més fort que els axiomes del llibre. Els capítols de Nivell~I no
en depenen, i el text ho assenyalarà cada vegada que calgui.
\end{nota}

Fins ara, tota família que hem format era interna perquè els seus valors
vivien dins un conjunt ja construït: parts d'un fix \(A\), elements d'un
exponencial, fibres d'un total donat. Els capítols següents arribaran, però,
a construccions on aquesta comoditat desapareix: una recursió transfinita els
valors de la qual són conjunts \emph{nous} sense codomini comú, el pas d'un
bon ordre al seu representant ordinal, una successió \(\alpha\mapsto\aleph_\alpha\).
En tots aquests casos tindrem una \emph{regla externa} ben definida
\(i\mapsto X_i\) i voldrem reunir-la en una \emph{família interna} \(E\to I\),
però cap axioma de Nivell~I no ens ho garantirà. El principi que ho permet,
que designarem provisionalment per \textbf{SR} (\enquote{substitució
estructural}), es pot enunciar orientativament així.

\begin{observacio}[Esquema de reunió de famílies (SR), provisional]\label{obs:sr}
Suposem que una fórmula, per a cada element \(i\) d'un conjunt \(I\),
determina un conjunt \(X_i\) \emph{llevat d'un únic isomorfisme compatible
amb les dades}. Aleshores ---diu SR--- existeix una aplicació \(p\colon E\to
I\) la fibra de la qual sobre cada \(i\) representa \(X_i\). És a dir: tota
regla externa uniformement i invariantment definida es pot representar per
una família interna.

No prenem això com a axioma en aquesta forma. L'enunciat definitiu, com un
\emph{esquema} de fórmules amb la condició de funcionalitat a isomorfisme
únic i l'estructura que la garanteix explicitades, és el
principi~\ref{sr:esquema} del capítol~\ref{cap:alephs}, que també n'identifica
la formulació adoptada a la literatura; aquí només el deixem anunciat
perquè el lector sàpiga on i per què farà falta.
\end{observacio}

\begin{observacio}[SR no és l'elecció]\label{obs:sr-no-eleccio}
Convé, ja d'entrada, no confondre SR amb l'axioma d'elecció, perquè fan
feines lògiques diferents ---la mateixa distinció que ja vam marcar al
capítol~\ref{cap:infinitud}
(observació~\ref{obs:eleccio-global-vs-numerable}). L'elecció tria un element
dins cada fibra d'una família que \emph{ja existeix} \(p\colon E\to I\)
(una secció, secció~\ref{sec:seccions}); SR, en canvi, garanteix
l'\emph{existència} de la família \(E\to I\) a partir d'una regla que defineix
els seus valors. La primera opera sobre un total donat; la segona el
construeix. En els capítols que vénen, mantindrem separats els dos usos i,
seguint la disciplina de la caixa d'axiomes, marcarem \enquote{elecció} i
\enquote{SR} allà on cadascun intervingui, mai indistintament.
\end{observacio}

\begin{resumcap}
\apartat{Idees centrals}
\begin{enumerate}
  \item Una família de conjunts indexada per \(I\) és una aplicació \(p\colon E\to I\); les fibres \(E_i=p^{-1}(i)\) són parts del total. La família constant és una projecció.
  \item Una secció de \(p\) és una funció d'elecció; per a fibres no buides, l'existència de seccions és l'axioma d'elecció. Quan les fibres tenen un element canònic (mínim, ínfim), la secció es construeix sense elecció.
  \item La suma d'una família és el total \(E\); existeix sempre. El producte és la part de \(E^I\) de les seccions: existeix sempre com a objecte, i té punts exactament quan \(p\) té secció.
  \item Una família de parts d'un conjunt fix és \(\alpha\colon I\to P(A)\); la unió i la intersecció indexades són \(\exists_{\pi_A}(R)\) i \(\forall_{\pi_A}(R)\), suprem i ínfim a \((P(A),\subseteq)\), amb les lleis de De Morgan; tot acotat dins \(A\).
  \item Una família de parts no buides, disjuntes i que cobreixen \(A\) és una partició (un epimorfisme \(A\twoheadrightarrow I\)), i recíprocament.
  \item Família interna (una \(p\colon E\to I\) donada) contra regla externa (una prescripció \(i\mapsto X_i\)): la regla de treball exigeix dir sempre quina aplicació representa la família. Quan els valors no tenen ambient comú, cal SR, que no és l'elecció: l'elecció tria dins fibres d'una família que existeix, SR fa existir la família.
\end{enumerate}
\apartat{Errors típics} Pensar que \enquote{una família de conjunts} és una regla i no una aplicació; confondre l'existència del producte amb la seva no-vacuïtat; llegir la unió indexada sobre elements generalitzats amb un \enquote{existeix} en el mateix estadi; escriure \(\bigcup_iX_i\) sense haver exhibit un total o un ambient; i confondre SR (existència de la família) amb l'elecció (secció d'una família existent).
\apartat{Lectura recomanada} \cite[cap.~8 i §§9.1--9.2]{lawvere-rosebrugh} per a les famílies com a aplicacions, les fibres i el pas als conjunts variables; \cite[§§9--10]{halmos} per a les famílies, les unions i els productes indexats en la teoria clàssica, on la distinció entre regla i família queda amagada per l'axioma de reemplaçament; \cite[§§2--3]{shulman-comparing} per a la formulació estructural de les famílies i del reemplaçament; i \cite[cap.~I]{maclane-moerdijk} per a les categories de fletxes \(\mathbf{Set}/I\) com a marc general de les famílies indexades.
\end{resumcap}
